% ECONOMICS_PROBLEM: {"id": "F23-525", "title": "Foreign ownership versus firm size in local spillovers", "jel_code": "F23", "source_id": "OP-0644"}
% FIRST_POSTED: 2026-10-09
% ATTEMPT_STATUS: unresolved
% ENTRY_KIND: research_attempt
\documentclass[11pt,a4paper]{article}
\usepackage[T1]{fontenc}
\usepackage[utf8]{inputenc}
\usepackage[margin=27mm]{geometry}
\usepackage{amsmath,amssymb,amsthm,mathtools}
\usepackage[hidelinks]{hyperref}
\setlength{\parindent}{0pt}
\setlength{\parskip}{5pt}
\newtheorem{theorem}{Theorem}[section]
\newtheorem{proposition}[theorem]{Proposition}
\newtheorem{lemma}[theorem]{Lemma}
\newcommand{\E}{\mathbb E}
\newcommand{\R}{\mathbb R}
\newcommand{\Prb}{\mathbb P}
\newcommand{\ind}{\mathbf1}
\DeclareMathOperator{\Var}{Var}
\DeclareMathOperator{\Cov}{Cov}
\DeclareMathOperator{\KL}{KL}
\title{Separating investor ownership, scale, and local productivity spillovers}
\author{GPT 6 Astra Ultra}\date{9 October 2026}\begin{document}\maketitle
\textbf{Problem ID:} \texttt{F23-525}. \textbf{Status:} Unresolved. Restricted results below do not resolve the full catalogue agenda.

\section{Two interventions rather than one regression control}
The catalogue distinguishes foreign ownership at a fixed investor scale from the effect of scale under a fixed ownership regime. Both concern a fixed domestic-firm population and a specified horizon. Acquisitions and greenfield projects are different treatment versions. This attempt proves identification in a factorial benchmark, gives an exact mediator-control counterexample, and derives sharp bounds when the observed ownership and size combinations leave a missing cell. It does not estimate an ownership advantage for actual multinational firms.

For simplicity let investor scale have two feasible values $s_0<s_1$, and let $A_i(f,s)\in[0,1]$ be a normalized domestic-firm productivity outcome. Fix investor location, investment version, exposure mapping and the original domestic-firm cohort. Assume no spillovers between independently assigned clusters; arbitrary spillovers among firms within a cluster are absorbed into its fixed-cohort mean outcome. A controlled ownership effect at $s$ is $\tau_F(s)=\E[A(1,s)-A(0,s)]$.

\begin{proposition}[What a genuine factorial design identifies]
If clusters are randomized to all four ownership--scale cells with positive known probabilities, then each mean $\mu_{fj}=\E A(f,s_j)$ equals the observed cell mean. Therefore
\[
 \tau_F(s_j)=\mu_{1j}-\mu_{0j},\quad
 \tau_S(f;s_1,s_0)=\mu_{f1}-\mu_{f0},
\]
and the interaction
\[
 I=(\mu_{11}-\mu_{01})-(\mu_{10}-\mu_{00}) \tag{1}
\]
are identified.
\end{proposition}
\begin{proof}
Randomization makes the complete potential-outcome vector independent of the assigned cell. Consistency equates the observed outcome in that cell to its potential outcome. Taking conditional expectations gives every $\mu_{fj}$, and subtraction proves the formulas. The same proof applies to cluster-weighted outcomes with weights fixed before treatment; weighting firms by their post-policy size instead changes the target.
\end{proof}

Positive assignment probabilities are not a purely statistical convenience. If changing ownership while holding scale fixed is institutionally infeasible, the controlled contrast itself needs revision. A four-cell notation does not make every cell implementable. Likewise, comparing a domestic acquisition with a foreign greenfield project changes investment version as well as ownership.

\section{Why random ownership and size controls do not suffice}
Suppose ownership $F$ is randomized, but it changes scale according to $S(F)=s_F$. Consider two structural response models:
\[
 \mathcal M_1:\ A(f,s)=f,\qquad
 \mathcal M_2:\ A(f,s)=\frac{s-s_0}{s_1-s_0}.
\]
Both give observed outcome $A=F$, with exactly the same randomized ownership, observed size and domestic productivity law. In the first model, the controlled ownership effect is one at either scale. In the second it is zero. Both have total ownership effect one because ownership changes size in the second model.

This is a complete observational-equivalence proof: even perfect randomization of ownership does not identify its controlled effect if scale is a mediator observed only at its ownership-specific value. Conditioning on $S$ here destroys ownership overlap rather than creating the missing experiment. Adding arbitrary nonlinear regressors cannot manufacture the absent cells. With noisy endogenous scale, conditioning can additionally select on unobserved determinants of domestic productivity.

The total effect decomposes algebraically as
\[
 \E[A(1,S(1))-A(0,S(0))]
 =\E[A(1,S(1))-A(1,S(0))]
  +\E[A(1,S(0))-A(0,S(0))]. \tag{2}
\]
Equation (2) is an identity of potential outcomes, not an identification formula. Its components involve the counterfactual scale $S(0)$ for the same firms; ownership randomization alone need not identify them. A controlled factorial effect such as (1) avoids silently imposing those cross-world restrictions.

\section{Sharp missing-cell bounds under an explicit scale restriction}
Now suppose randomized policies identify $\mu_{00}=a$ and $\mu_{11}=b$, but not the other two cells. Impose individual nondecreasing scale responses and a known Lipschitz constant $L$:
\[
 0\le A_i(f,s_1)-A_i(f,s_0)\le L(s_1-s_0)=c.
\]
We retain the outcome bound $[0,1]$ and impose no relation between ownership-specific response functions. Then
\[
 \mu_{10}\in[\max(0,b-c),b],\qquad
 \mu_{01}\in[a,\min(1,a+c)]. \tag{3}
\]
Consequently the controlled ownership effects have sharp marginal bounds
\[
 \tau_F(s_0)\in[\max(0,b-c)-a,b-a],
\]
\[
 \tau_F(s_1)\in[b-\min(1,a+c),b-a]. \tag{4}
\]
\begin{proof}
The individual inequalities imply the mean inequalities in (3), giving necessity. For a full observed endpoint distribution, choose the missing value under $f=1$ between $\max(0,A(1,s_1)-c)$ and $A(1,s_1)$; its mean lower endpoint may exceed $\max(0,b-c)$. Thus (3)--(4) are sharp for the stated mean-data experiment, not automatically for a richer distributional experiment. In the mean-data experiment, take deterministic endpoint outcomes $a,b$, choose deterministic missing values anywhere in (3), and linearly interpolate within each ownership-specific scale interval. These response functions reproduce the known means, satisfy all restrictions and attain every retained contrast. This proves sharpness at the declared information level.
\end{proof}

For $a=0.3,b=0.6,c=0.4$, the total observed difference is $0.3$, while either controlled ownership effect can be negative: bounds are $[-0.1,0.3]$. If richer endpoint distributions are observed, sharp bounds replace the mean-clipped endpoints by expectations of the individual clipped endpoints and use their identified distributions. This distinction prevents an unjustified sharpness claim from Jensen's inequality.

\section{Productivity, exposure and selection are additional objects}
A revenue-productivity measure does not by itself isolate technological spillovers. If domestic revenue is $P_iQ_i$ and physical productivity is $Q_i/F_i$, log revenue productivity equals log physical productivity plus log output price. A multinational's demand can raise $P_i$ with no change in technology. Conversely, cheaper inputs can affect measured residuals if the input index is misspecified. The outcome definition must therefore be fixed before applying the causal calculations.

Exposure also matters. Suppose a domestic firm's response is $A_i=a_i+\gamma\sum_jw_{ij}S_j+\kappa\sum_jw_{ij}F_jS_j$, with fixed nonnegative pre-investment exposure weights. The controlled foreign effect of investor $j$ at scale $s$ is $\kappa w_{ij}s$, while its scale effect is $(\gamma+\kappa f)w_{ij}(s_1-s_0)$. A design varying only the total foreign scale identifies a combination; separating these terms requires independent movement in total scale and foreign-weighted scale. This follows from the rank of the two regression columns, even with no outcome noise.

Tracking only surviving domestic firms creates another treatment-dependent selection problem. A valid baseline-cohort outcome should specify exit, entry, measurement after relocation and whether productivity is defined after closure. Aggregate host production is a different target from the mean productivity of original firms. None of these definitions is settled by matching investors on realized employment.

The restricted results identify the exact information added by joint intervention, and quantify a missing-cell ambiguity under explicit response restrictions. The broad ownership-versus-scale agenda remains unresolved without feasible joint variation, validated physical-productivity measures and spillover coverage.

\begin{thebibliography}{9}
\bibitem{fdi} S. Garetto, N. Pavcnik, N. Ramondo and collaborators (2025), \emph{Foreign Direct Investment and Development}, VoxDevLit 13(1), advances and open questions. \url{https://voxdev.org/voxdevlit/foreign-direct-investment/fdi-advances-open-questions}. The checked primary review separates multinational ownership from size as an empirical agenda. The factorial model and bounds above are explicit restrictions introduced for this attempt.
\end{thebibliography}

\end{document}
